\section{The coupled datum-defect classification}
\label{sec:theorem}

This section states and proves the paper's proven contribution: a \emph{basis-invariant}
classification of the null space of the coupled Fisher information \eqref{eq:fisher} into
its spatial, temporal, and genuinely coupled parts; the exact datum gauge of a fully
internal \emph{free-network} (un-anchored, mutual-clock) lunar problem under the SI
convention; the precise criterion that certifies a coupled defect; and the conditions under
which external ties lift each temporal defect---connecting the free-network gauge to the
tied, defect-$0$ realization of Section~\ref{sec:anchor}. Everything here is \Proven{}
algebra; every physical interpretation is labelled where it appears.

\subsection{The coupled datum defect and its projector}
\label{sec:defectdef}

\begin{definition}[Coupled datum defect]
\label{def:defect}
Let $\Nsp=\ker\Imat(\thetav)$ be the null space of the coupled Fisher information, of
dimension $d=\dim\Nsp$ (the \emph{datum defect}), and let $\Umat\in\R^{9\times d}$ have
orthonormal columns spanning $\Nsp$. With the spatial and temporal subspaces
$\Ssp,\Tsp$ of \eqref{eq:subspaces}, write $\Umat_{\Ssp}=\Umat[0{:}7,:]\in\R^{7\times d}$
and $\Umat_{\Tsp}=\Umat[7{:}9,:]\in\R^{2\times d}$ for the row-blocks of $\Umat$ on the two
subspaces. Define
\begin{align}
  \dim(\Nsp\cap\Ssp) &\;=\; d-\rank(\Umat_{\Tsp}) &&\text{(purely spatial defects)},
  \label{eq:dimspatial}\\
  \dim(\Nsp\cap\Tsp) &\;=\; d-\rank(\Umat_{\Ssp}) &&\text{(purely temporal defects)},
  \label{eq:dimtemporal}\\
  \operatorname{coupled\,dim} &\;=\; \rank(\Umat_{\Ssp})+\rank(\Umat_{\Tsp})-d
     &&\text{(coupled defects)}.
  \label{eq:coupleddim}
\end{align}
Equivalently, and manifestly basis-free, let $\Pnull=\Umat\Umat^\top$ be the orthogonal
projector onto $\Nsp$ and let $\Pst=\Pnull[0{:}7,\,7{:}9]\in\R^{7\times2}$ be its
spatial--temporal off-diagonal block. A coupled defect is \emph{certified} by
$\Pst\neq\bm0$; the magnitude $\lVert\Pst\rVert_{\mathrm F}$ is an orthogonally-invariant
(but, by Lemma~\ref{lem:precond}, not preconditioner-invariant) measure of its strength.
\end{definition}

The three counts in \eqref{eq:dimspatial}--\eqref{eq:coupleddim} partition the defect: a
purely spatial null direction is invisible in the temporal rows, a purely temporal one is
invisible in the spatial rows, and a coupled direction is visible in both. Their sum is
$\dim(\Nsp\cap\Ssp)+\dim(\Nsp\cap\Tsp)+\operatorname{coupled\,dim}=d$ \emph{identically},
by the definition of $\operatorname{coupled\,dim}$ in \eqref{eq:coupleddim}; no
``residual'' case exists. That count is moreover nonnegative: the vertical stack
$[\Umat_{\Ssp};\Umat_{\Tsp}]$ of the two row-blocks is $\Umat$ itself, so
$\rank[\Umat_{\Ssp};\Umat_{\Tsp}]=\rank\Umat=d\le\rank(\Umat_{\Ssp})+\rank(\Umat_{\Tsp})$
by rank subadditivity under row stacking, giving
$\operatorname{coupled\,dim}=\rank(\Umat_{\Ssp})+\rank(\Umat_{\Tsp})-d\ge0$. The projector
form is the object the engine reports, and it is the right object because of the following
invariance, which is the paper's central rigor point.

\subsection{Basis invariance}
\label{sec:invariance}

A null space does not come with a canonical basis: any $\Umat\bm Q$ with $\bm Q$ orthogonal
spans the same $\Nsp$, and an eigensolver returns one arbitrary such $\Umat$. A
classification that depended on the individual columns of $\Umat$ would therefore be
ill-posed. The naive quantity ``the energy of a null \emph{vector} in the temporal block,''
$\lVert\bm u_{\Tsp}\rVert$ for a chosen null vector $\bm u$, is exactly such an ill-posed
quantity---it changes under $\Umat\mapsto\Umat\bm Q$. The classification of
Definition~\ref{def:defect} avoids this trap.

\begin{proposition}[Basis invariance of the classification]
\label{prop:invariance}
The three dimensions \eqref{eq:dimspatial}--\eqref{eq:coupleddim} and the coupling strength
$\lVert\Pst\rVert_{\mathrm F}$ depend only on the subspace $\Nsp$, not on the choice of
orthonormal basis $\Umat$. Formally, they are invariant under $\Umat\mapsto\Umat\bm Q$ for
every orthogonal $\bm Q\in\R^{d\times d}$.
\end{proposition}

\begin{proof}
Let $\bm P_{\Ssp}$ and $\bm P_{\Tsp}$ be the fixed coordinate projectors onto $\Ssp$ and
$\Tsp$ (selecting rows $0$--$6$ and $7$--$8$). For the row-block $\Umat_{\Tsp}=\bm P_{\Tsp}\Umat$,
$\rank(\Umat_{\Tsp})=\dim(\bm P_{\Tsp}\Nsp)$, which is a property of the subspace $\Nsp$ and the
fixed projector, independent of the basis; the same holds for $\rank(\Umat_{\Ssp})$. Hence
\eqref{eq:dimspatial}--\eqref{eq:coupleddim}, being integer combinations of these ranks and
$d=\dim\Nsp$, are subspace invariants. For the projector, orthogonal $\bm Q$ gives
$(\Umat\bm Q)(\Umat\bm Q)^\top=\Umat\bm Q\bm Q^\top\Umat^\top=\Umat\Umat^\top=\Pnull$, so
$\Pnull$ and therefore its fixed sub-block $\Pst$ and $\lVert\Pst\rVert_{\mathrm F}$ are
literally unchanged. (Indeed $\Pnull$ is the unique orthogonal projector with range $\Nsp$,
so it is intrinsic to the subspace.)
\end{proof}

\begin{remark}[The projector zeroness certifies coupling; its norm measures strength]
\label{rem:projnorm}
$\lVert\Pst\rVert_{\mathrm F}=0$ if and only if $\Pnull$ is block-diagonal in
$\Ssp\oplus\Tsp$, i.e.\ $\Nsp=(\Nsp\cap\Ssp)\oplus(\Nsp\cap\Tsp)$ is a direct sum of a
spatial and a temporal part with no direction mixing the two. This \emph{binary} certificate
$\Pst=\bm0$ is the invariant object: it is unchanged under both orthogonal basis rotation
(Proposition~\ref{prop:invariance}) and the positive-diagonal preconditioner
(Lemma~\ref{lem:precond}). A single null vector $\bm u=(\bm u_{\Ssp};\bm u_{\Tsp})$ with both
blocks nonzero forces $\Pst\neq\bm0$ and contributes to $\operatorname{coupled\,dim}$. The
\emph{value} $\lVert\Pst\rVert_{\mathrm F}$ is a strength magnitude, invariant under
orthogonal basis change but not under the preconditioner; the engine reports
$\lVert\Pst\rVert_{\mathrm F}=0.5$ for the canonical single coupled direction
$\bm u=(\bm e_{\text{scale}}+\bm e_{\dtau})/\sqrt2$ on an \emph{un-preconditioned} synthetic
matrix, and a unit test confirms this value is unchanged when the null basis is rotated by an
arbitrary angle (Section~\ref{sec:anchor}). For the anchored networks $\Pst=\bm0$, so the
strength value enters no validated number.
\end{remark}

\subsection{The classification theorem and the exact defect count}
\label{sec:classify}

\begin{theorem}[Coupled free-network datum gauge: classification and count]
\label{thm:classify}
Consider a fully internal \emph{free-network} lunar PNT problem: mutual one-way
pseudoranges among orbiters and surface beacons with common-mode clocks and \emph{no}
external position or time tie (model~(A) of Section~\ref{sec:models}), under the SI
convention of Section~\ref{sec:si}, with coupled Fisher \eqref{eq:fisher} assembled from the
free/differential rows in which a common timescale offset and rate cancel \eqref{eq:free}.
Then:
\begin{enumerate}
\item[(i)] \textbf{Exact defect.} The datum gauge defect is
$d=6+2=8$: three translations, three rotations, the clock offset $\dtau$, and the clock
rate $\dalpha$. Scale $s$ is \emph{observable} (not a defect), because under SI a
pseudorange is scale-bearing through the fixed $c$.
\item[(ii)] \textbf{Temporal defects are pure.} In the free/differential design the clock
term enters each one-way range as a difference (receiver minus transmitter), so a common
offset $\dtau$ and a common rate $\dalpha$ cancel in every internal observation and their
datum columns are \emph{identically zero} \eqref{eq:free}; hence $\bm e_{\dtau}$ and
$\bm e_{\dalpha}$ lie in $\ker\Imat$ exactly, giving $\dim(\Nsp\cap\Tsp)=2$ and
$\operatorname{coupled\,dim}=0$ for the free network alone. (In the \emph{anchored} model~(B),
by contrast, the offset column is $+1$---an absolute offset against a known reference
\eqref{eq:oneway}---and is therefore observable; that model is treated in
Section~\ref{sec:anchor}.)
\item[(iii)] \textbf{The joint gauge is not a direct sum in general.} Once an external tie
that couples the subspaces is admitted---in particular the relativistic rate tie
\eqref{eq:ratetie}, whose scale column is nonzero---the null space acquires a direction with
$\Pst\neq\bm0$, so the joint frame$\oplus$timescale gauge is \emph{not} the direct sum of
the spatial gauge of \citep{baweja2026datum} and the temporal gauge.
\end{enumerate}
\end{theorem}

\begin{proof}
\emph{(i), (ii).} Under SI (Section~\ref{sec:si}) the scale column of every one-way row is
the physical range's dependence on the lunar radius through the fixed $c$, which is nonzero
and geometry-dependent across a network of distinct line-of-sight directions; hence scale
is not annihilated and contributes no null direction (this is the P1 result that a
geometrically diverse network observes scale \citep{baweja2026datum}, imported unchanged).
The rigid-body translation ($3$) and rotation ($3$) directions are the classical Helmert
datum defect of a free network: a common translation and a common infinitesimal rotation of
the point set are unobservable to internal ranges \citep{baarda1973,teunissen2000}, giving
six spatial null directions. For the temporal columns, the free/differential design
\eqref{eq:free} carries the clock term as a receiver-minus-transmitter difference, so a
common offset $\dtau$ and a common rate $\dalpha$ leave every internal observation
invariant---their datum columns are \emph{identically zero}---and $\bm e_7=\bm e_{\dtau}$,
$\bm e_8=\bm e_{\dalpha}$ are annihilated by $\Imat$. (This is the defining property of the
free network; it is \emph{not} the anchored row \eqref{eq:oneway}, whose $+1$ offset column
makes the offset observable, since there $(\Imat\bm e_7)_7=\sum_i g_{7,i}^2=N>0$.) Thus
$\ker\Imat\supseteq\operatorname{span}\{\text{6 rigid},\bm e_7,\bm e_8\}$, $d\ge8$; and
because a geometrically diverse free network observes everything else (scale and all
relative geometry), $d=8$ with the temporal pair lying entirely in $\Tsp$, so
$\dim(\Nsp\cap\Tsp)=2$ and $\operatorname{coupled\,dim}=0$.

\emph{(iii).} Add the rate tie \eqref{eq:ratetie}. Its row
$\bm g^{\mathrm{rt}}=-\tfrac{\partial\alpha}{\partial s}\,\bm e_{s}+\bm e_{\dalpha}$ has a
nonzero scale component, so it constrains the combination
$\dalpha-\tfrac{\partial\alpha}{\partial s}\,\delta s$. If scale is otherwise fixed, this
removes the $\dalpha$ defect (below); in the free-network sub-case where no independent
scale tie is present, the surviving null direction is
$\bm u\propto\bigl(\tfrac{\partial\alpha}{\partial s}\,\bm e_{s}+\bm e_{\dalpha}\bigr)$,
which has a nonzero spatial block ($\bm e_s\in\Ssp$) and a nonzero temporal block
($\bm e_{\dalpha}\in\Tsp$); by Remark~\ref{rem:projnorm}, $\Pst\neq\bm0$. Hence the joint
gauge contains a direction mixing $\Ssp$ and $\Tsp$ and is not $\Ssp$-$\Tsp$ block-diagonal.
\end{proof}

\begin{remark}[The off-diagonal criterion: necessary, not sufficient]
\label{rem:necessary}
It is tempting to certify a coupled defect from the \emph{Fisher} rather than its null
space---by checking whether the Fisher's spatial--temporal block $\Imat_{\Ssp\Tsp}$ is
nonzero. This is wrong. $\Imat_{\Ssp\Tsp}\neq\bm0$ (the Fisher is not block-diagonal) is
\emph{necessary} for a coupled defect but \emph{not sufficient}: the coupling can live
entirely among the \emph{observable} directions, leaving the null space a clean direct sum.
Concretely, a network whose scale$\leftrightarrow$offset block is nonzero (they are
correlated) but full-rank has $\Imat_{\Ssp\Tsp}\neq\bm0$ and yet $\operatorname{coupled\,dim}=0$,
$\Pst=\bm0$---the very case of the well-posed anchor network of
Section~\ref{sec:anchor}, which has $\operatorname{coupled\,dim}=0$ despite a
non-block-diagonal Fisher. A coupled \emph{defect} is certified \emph{only} through the null
space: by $\Pst\neq\bm0$, equivalently $\operatorname{coupled\,dim}>0$ via the rank formula
\eqref{eq:coupleddim}.
\end{remark}

\subsection{Lifting the temporal defects}
\label{sec:lift}

The temporal defects are lifted by external ties, and the mechanism is the intersection law
for the null spaces of summed positive-semidefinite (PSD) information blocks.

\begin{lemma}[Intersection law]
\label{lem:intersect}
For PSD matrices $\Imat_A,\Imat_B\succeq0$,
$\ker(\Imat_A+\Imat_B)=\ker\Imat_A\cap\ker\Imat_B$.
\end{lemma}

\begin{proof}
For PSD $\Imat_A,\Imat_B$ and any $\bm v$, $\bm v^\top(\Imat_A+\Imat_B)\bm v=
\bm v^\top\Imat_A\bm v+\bm v^\top\Imat_B\bm v$ is a sum of two nonnegative terms, so it
vanishes iff both vanish, i.e.\ iff $\Imat_A\bm v=\bm0$ and $\Imat_B\bm v=\bm0$ (a PSD matrix
annihilates exactly the vectors of zero quadratic form). Hence
$\ker(\Imat_A+\Imat_B)=\ker\Imat_A\cap\ker\Imat_B$.
\end{proof}

\begin{theorem}[Lift of the free-network temporal defects]
\label{thm:lift}
Let $\Imat$ be the \emph{free-network} Fisher of Theorem~\ref{thm:classify}, whose temporal
gauge is $\ker\Imat\supseteq\operatorname{span}\{\bm e_{\dtau},\bm e_{\dalpha}\}$ (the two
free-network temporal defects). External ties reduce this gauge, connecting the free-network
defect of eight to the fully-tied, defect-$0$ realization of Section~\ref{sec:anchor}:
\begin{enumerate}
\item[(a)] \textbf{Offset.} Adding a time-tie block \eqref{eq:timetie} (rank-one at
$\bm e_{\dtau}$) removes $\bm e_{\dtau}$ from the null space: by Lemma~\ref{lem:intersect} the
combined null space is $\ker\Imat\cap\ker\Imat_{\mathrm{tt}}$, and $\Imat_{\mathrm{tt}}$
annihilates every vector except those with a $\dtau$ component, so $\bm e_{\dtau}$ is lifted
while $\bm e_{\dalpha}$ survives.
\item[(b)] \textbf{Rate.} Adding the rate tie \eqref{eq:ratetie} constrains
$\dalpha-\tfrac{\partial\alpha}{\partial s}\,\delta s$. Consider the state \emph{after} the
offset has been tied (part~a) and with scale observable (as under SI) or independently tied:
the surviving null space then has no scale component ($\delta s=0$) and its only temporal
direction is $\bm e_{\dalpha}$, so the rate tie forces $\dalpha=0$ and lifts it fully. Were
scale instead left free---a free-network sub-case, not the SI default---the rate tie alone
would fuse the rate and scale directions into a single coupled scale$\leftrightarrow$rate
null direction (reducing that sector's dimension by one) rather than remove the defect. Thus
(a) followed by (b) lifts \emph{both} free-network temporal defects, and the fully-tied
network reaches defect $0$ (Section~\ref{sec:anchor}).
\end{enumerate}
\end{theorem}

\begin{proof}
Both parts are Lemma~\ref{lem:intersect} applied to $\Imat+\Imat_{\mathrm{tie}}$. For (a),
$\Imat_{\mathrm{tt}}=\bm e_{\dtau}\bm e_{\dtau}^\top$ has kernel
$\{\bm v:v_{\dtau}=0\}$; intersecting with $\ker\Imat$ removes exactly $\bm e_{\dtau}$. For
(b), $\Imat_{\mathrm{rt}}=\bm g^{\mathrm{rt}}(\bm g^{\mathrm{rt}})^\top$ has kernel
$\{\bm v:\bm g^{\mathrm{rt}}\!\cdot\!\bm v=0\}=\{\bm v:v_{\dalpha}=\tfrac{\partial\alpha}
{\partial s}v_{s}\}$; intersecting with $\ker\Imat$, and using that under SI $\ker\Imat$ has
no scale component ($v_s=0$ on the null space), forces $v_{\dalpha}=0$, removing
$\bm e_{\dalpha}$. Without an independent scale tie, $v_s$ is free and the surviving
direction is $\tfrac{\partial\alpha}{\partial s}\bm e_s+\bm e_{\dalpha}$, a single coupled
direction.
\end{proof}

\begin{remark}[Exact-null lift versus near-defect lift]
\label{rem:distinct}
The lifts of Theorem~\ref{thm:lift} act on \emph{exact} null directions (defects). They are
distinct from, and must not be conflated with, the raising of a small-but-positive
eigenvalue of the $\{s,\dtau\}$ marginal by additional observations---the
\emph{near}-defect of coupling~1 (Section~\ref{sec:coupling1}), which is never a member of
$\ker\Imat$ in a diverse network. One is a change of the defect $d$; the other is a change
of a strictly positive eigenvalue.
\end{remark}
